\documentclass[12pt]{article}
\usepackage[utf8]{inputenc}
\usepackage[T1]{fontenc}
\usepackage{hyperref}
\title{Xer\'ente Language Resources for K-12 Curriculum Development}
\author{Compiled by Claude Sonnet 5.5 (claude-sonnet-5-5) for Andr\'as Kornai}
\date{}
\begin{document}
\maketitle

\section{Introduction}
Xer\'ente (Akw\~e; ISO 639-3 xer; Glottocode xere1240) is a J\^e language of the Xerente Indigenous Area in Tocantins, Brazil. One source in the search results says the population is about 3,100 people. The language is vigorous and spoken by the whole community, and it is taught in some schools. Most of the resources found are in Portuguese.

\section{Primary Educational Resources}
\begin{itemize}
\item \emph{Dicion\'ario Escolar Xerente--Portugu\^es / Portugu\^es--Xerente}, a school dictionary. It is known from a 2013 \emph{Entreletras} (Aragua\'ina, TO) paper analysing its macro- and microstructure: \url{https://sumarios.org/node/145954}
\item Rinaldo de Mattos, \emph{L\'ingua Xerente: como se l\^e e como se escreve}, a mother-tongue grammar and reading guide made with Xerente teachers and students. This was mentioned in search summaries, but no URL for the work itself was seen, so it is not in the JSON.
\end{itemize}

\section{Linguistic and Reference Materials}
\begin{itemize}
\item Sinval Martins de Sousa Filho, \emph{Aspectos morfossint\'aticos da l\'ingua Akw\~e-Xerente (J\^e)} (2010, from a 2007 UFG doctoral thesis). No URL was seen for it.
\item A 2018 study proposing a standard spelling for nouns in Akw\~e Mrm\~ez\'e and the lack of a standard among teachers. The search listed several journal PDFs, but I could not match this study to a specific URL, so it is not in the JSON.
\end{itemize}

\section{Religious and Community Resources}
None were found in the search results.

\section{Digital and Audio Resources}
No audio, video or software resources were found.

\section{Curriculum Development Context}
Xer\'ente is vigorous and already taught as a subject in some schools. Brazilian indigenous school education (educa\c{c}\~ao escolar ind\'igena) supports bilingual instruction. A curriculum could build on the existing school dictionary and on the mother-tongue reading and writing guide, with Xerente teachers. Spelling is not yet standardised among teachers, so an orthography agreement should come first.
\end{document}